\subsubsection{One fixed continuum operator controls every finite grid}
\begin{lemma}[Uniform Cauchy-cell comparison]\label{deg:lem:cell-grid}
Fix $0<\alpha<1$, let $a=\sqrt2$, and let $S$ be the positive integral
operator on $L^2(0,a)$ with kernel $\sqrt{xy}/(x+y)$.
Suppose $x_1,\ldots,x_N\in(0,2]$ and $h>0$ satisfy $hN\le1$.
For all low nodes $x_i\le a$, suppose $x_i\ge2h$ and the intervals
$I_i=[x_i-h,x_i]$ are pairwise disjoint. Put
\[
 T_{\alpha,h}(i,j)=hK_\alpha(\log x_i^2-\log x_j^2),\qquad
 M_\alpha=\max\{\sin(\pi\alpha),\tan(\pi\alpha/2)\}.
\]
Uniformly over all such finite grids,
\[
 \N_{T_{\alpha,h}}(\varepsilon)
 \le\N_S\!\left(\frac{\varepsilon}{4M_\alpha\log2}\right)
       +O_\alpha(1+\log^+(1/\varepsilon)),
 \qquad \log^+t=\max(0,\log t).
\]
Consequently their uniform upper count coefficient is at most
$1/(2\pi^2)$ at scale $\log^2(1/\varepsilon)$.
\end{lemma}
\begin{proof}
The ratio of the Fourier symbol in~\eqref{deg:eq:symbol} to its value
at $\alpha=1/2$ is
\[
 \frac{\sin(\pi\alpha)\cosh(2\pi\xi)}
      {\cosh(2\pi\xi)+\cos(\pi\alpha)}\le M_\alpha.
\]
Fourier inversion gives, on every finite real grid $t_i$,
\[
 \sum_{i,j}\overline v_i v_jK_\alpha(t_i-t_j)
  =\frac1{2\pi}\int_{\R}\widehat K_\alpha(\xi)
                   \left|\sum_jv_je^{i\xi t_j}\right|^2d\xi.
\]
Thus $T_{\alpha,h}\preceq M_\alpha T_{1/2,h}$ in Loewner order,
including for atomic grids. It suffices to consider $\alpha=1/2$, when
\[
 T_{1/2,h}(i,j)=h\frac{\sqrt{x_ix_j}}{x_i+x_j}.
\]

For each low node define
\[
 f_i(x)=\frac{\sqrt{x_i}}{\sqrt h\sqrt x}\,1_{I_i}(x).
\]
These functions are orthogonal, and
\[
 d_i^2=\norm{f_i}^2
   =\frac{x_i}{h}\log\frac{x_i}{x_i-h}\le2\log2.
\]
Indeed $r\log(r/(r-1))$ decreases for $r>1$, and $x_i/h\ge2$.
The Laplace formula for $(x+y)^{-1}$ shows that $S$ is positive and
trace class, with trace $a/2$. Its compressed Gram matrix is
\[
 B_{ij}=\langle f_i,Sf_j\rangle
 =h\sqrt{x_ix_j}\int_0^\infty e^{-t(x_i+x_j)}\phi(ht)^2dt,
 \qquad \phi(z)=\frac{e^z-1}{z},\quad\phi(0)=1.
\]
Every cell has the same length, so the multiplier is common to all
coordinates. Since $\phi(z)\ge1$ for $z\ge0$,
\[
 v^*(B-T_{\rm low})v
 =h\int_0^\infty[\phi(ht)^2-1]
       \left|\sum_i v_i\sqrt{x_i}e^{-tx_i}\right|^2dt\ge0.
\]
The integrals converge because $x_i\ge2h$. This is a positive
quadratic-form comparison, rather than an entrywise inference.
Normalize $f_i/d_i$ to obtain an isometry $J$ and write
$B=D J^*SJ D$, $D=\diag(d_i)$.
Compression min--max and the singular values of $S^{1/2}JD$ give
\[
 \lambda_j(T_{\rm low})\le\lambda_j(B)
                \le2\log2\,\lambda_j(S).
\]
No Loewner inequality $DTD\preceq\norm D^2T$ is used.

On the high block $a<x_i\le2$, expand
\[
 \frac1{x+y}=\sum_{\ell=0}^{q-1}
       \frac{(-1)^\ell(x-a)^\ell}{(a+y)^{\ell+1}}+R_q(x,y),
 \qquad |R_q(x,y)|\le\frac{r_0^q}{2a},\quad
 r_0=\frac{2-a}{2a}<1.
\]
After multiplication by $h\sqrt{xy}$ this is a rank-$q$ matrix
approximation whose entry errors are at most $hr_0^q/a$.
Since the block has at most $N$ rows and $hN\le1$, its operator error
is at most $r_0^q/a$. The best-rank singular-value bound, also valid
for this possibly nonsymmetric approximant, gives
\[
 \N_{T_{\rm high}}(t)=O(1+\log^+(1/t)).
\]
Positive pinching, Lemma~\ref{deg:lem:pinch}, bounds the full count by
the sum of the two block counts at half the threshold. This proves the
displayed estimate, first at $\alpha=1/2$ and then for every fixed
$\alpha$ by the Fourier comparison.

Laptev's classical weighted-Cauchy theorem~\cite[Theorem (3)]{Laptev}
with exponent one and weights $x^{1/2}y^{1/2}$ has parameter
$\rho=1$ and gives
\[
 \N_S(\varepsilon)\sim\frac1{2\pi^2}\log^2(1/\varepsilon).
\]
Rescaling $(0,1)$ to $(0,a)$ only multiplies eigenvalues by a fixed
factor. The symmetric exponent-one case is explicitly attributed by
Laptev to Widom's earlier Hankel theorem. Crucially, $S$ is fixed before
taking the supremum over grids. Constant threshold factors and the
uniform logarithmic remainder do not change its leading coefficient.
\end{proof}

\begin{proof}[Sharp uniform upper in Theorem~\ref{deg:thm:main}]
For a path on $n$ vertices set
\[
 h=1/n,\qquad x_k=2\sin(k\pi/(2n)),\quad1\le k<n.
\]
The low nodes have $k\le\lfloor n/2\rfloor$, $x_1\ge2h$, and
successive spacing at least $\pi/(\sqrt2n)>2h$.
Their left cells satisfy Lemma~\ref{deg:lem:cell-grid}; also
$h(n-1)\le1$. The derivative in~\eqref{deg:eq:kernel} is
$R T_{\alpha,h}R$, with $R_{kk}=\sqrt{nz_k}$ and $\norm R^2\le2$.
The sandwich~\eqref{deg:eq:sandwich} and bounded-congruence eigenvalue
comparison give $\lambda_j(C)\le2\lambda_j(T_{\alpha,h})$.
Hence the upper coefficient is at most $1/(2\pi^2)$, uniformly over
all path sizes and all edited locations.

For a cycle on $n\ge3$ vertices, combine each equal eigenvalue pair
$\lambda_k=\lambda_{n-k}$ into its one active incidence coordinate.
Its leverage weight is $2/n$; if $n$ is even, the singleton
$\lambda_{n/2}=4$ has weight $1/n$.
Set
\[
 h=2/n,\qquad x_k=2\sin(\pi k/n),\quad
                       1\le k\le\lfloor n/2\rfloor.
\]
Now $h\lfloor n/2\rfloor\le1$. Low nodes have
$k\le\lfloor n/4\rfloor$, $x_1\ge4/n=2h$ whenever this set is nonempty,
and spacing at least $\sqrt2\pi/n>2h$.
Thus the same lemma applies, including when the low set is empty.
The active derivative is a diagonal congruence of $T_{\alpha,h}$
with norm at most one; its even endpoint factor is $1/\sqrt2$.
All other incidence coordinates have zero response.
The same uniform upper coefficient therefore holds over all cycles.
Every connected simple graph of maximum degree two with an edge is a
path or cycle. The centered-path lower bound below matches this uniform
upper for both $W_{\mathrm{path},\alpha}$ and $W_{2,\alpha}$.
\end{proof}

\begin{remark}
A nonsharp logarithmic-squared upper can alternatively be obtained from
low-energy leverage truncation, Stahl's root-exponential rational
approximation~\cite{Stahl}, and the classical rational update rank
identity discussed in~\cite[Section 3.3]{BCKS}.
The equal-width cell argument proves the sharp coefficient without
assuming a uniform finite-grid Nystr\"om convergence rate.
\end{remark}
